\Askhsh{14406}{4}
\begin{ylh}{2}
    \item 2.1. Οι Πράξεις και οι Ιδιότητές τους
    \item 3.1. Εξισώσεις 1ου Βαθμού
    \item 3.3. Εξισώσεις 2ου Βαθμού
\end{ylh}
Δίνονται οι μη μηδενικοί πραγματικοί αριθμοί α, β, με $a \neq \beta$ για τους οποίους ισχύει: $\dfrac{a^{2} + 1}{\beta^{2} + 1} = \dfrac{a}{\beta}$ .
\begin{alist}
\item Να αποδείξετε ότι οι αριθμοί α και β είναι αντίστροφοι.
\monades{5}
\item Να υπολογίσετε την τιμή της παράστασης Κ = $\dfrac{a^{22} \cdot \left( \beta^{3} \right)^{8}}{a^{- 2} \cdot (a\beta)^{25}}$.
\monades{7}
\item Αν επιπλέον οι μη μηδενικοί αριθμοί α και β εκφράζουν τα μήκη των πλευρών ορθογωνίου παραλληλογράμμου με άθροισμα $\dfrac{5}{2}$, να τους υπολογίσετε.
\monades{8}
\item Να βρείτε τον αριθμό που πρέπει να προσθέσετε στο α ή στο β, έτσι ώστε το ορθογώνιο παραλληλόγραμμο να γίνει τετράγωνο.
\monades{5}
\end{alist}

